\ifdefined\gkver\else\def\gkver{student}\fi
\documentclass[\gkver]{gaokaozhenti}
\newcommand{\ccnum}[1]{{\fontspec{Noto Sans CJK JP}#1}}
\begin{document}

\chapter{2020年全国理综Ⅱ}
%% paperType: 理综物理部分

\begin{enumerate}
\item
%% number: 1
%% paperName: 2020年全国理综Ⅱ第 1 题 6 分
%% typeId: 1
%% type: 单选题
%% chapter: 电磁感应
%% point: 电磁感应现象
%% method: 无
%% score: 6
%% degree: 950
%% duplicateId: 0
%% body:
管道高频焊机可以对由钢板卷成的圆管的接缝实施焊接。焊机的原理如图所示，圆管通过一个接有高频交流电源的线圈，线圈所产生的交变磁场使圆管中产生交变电流，电流产生的热量使接缝处的材料熔化将其焊接。焊接过程中所利用的电磁学规律的发现者为\xzanswer{D}

\onepicture[w=6.50cm]{figs/01.png}

\fourchoices[answer=D]
{库仑}
{霍尔}
{洛伦兹}
{法拉第}
%% answer: D
%% memo:
\memoanswer{
由题意可知，圆管为金属导体，导体内部自成闭合回路，且有电阻，当周围的线圈中产生出交变磁场时，就会在导体内部感应出涡电流，电流通过电阻要发热。该过程利用的原理是电磁感应现象，其发现者为法拉第。故选 D。
}

\item
%% number: 2
%% paperName: 2020年全国理综Ⅱ第 2 题 6 分
%% typeId: 1
%% type: 单选题
%% chapter: 曲线运动
%% point: 人造地球卫星
%% method: 无
%% score: 6
%% degree: 650
%% duplicateId: 0
%% body:
若一均匀球形星体的密度为 $\rho$，引力常量为 $G$，则在该星体表面附近沿圆轨道绕其运动的卫星的周期是\xzanswer{A}

\fourchoices[answer=A]
{$\sqrt{\frac{3\pi}{G\rho}}$}
{$\sqrt{\frac{4\pi}{G\rho}}$}
{$\sqrt{\frac{1}{3\pi G\rho}}$}
{$\sqrt{\frac{1}{4\pi G\rho}}$}
%% answer: A
%% memo:
\memoanswer{
卫星在星体表面附近绕其做圆周运动，由万有引力提供向心力，有

$\frac{GMm}{R^{2}}=m\frac{4\pi^{2}}{T^{2}}R$

又 $M=\rho V=\rho\cdot\frac{4}{3}\pi R^{3}$，代入上式解得

$T=\sqrt{\frac{3\pi}{G\rho}}$

故选 A。
}

\item
%% number: 3
%% paperName: 2020年全国理综Ⅱ第 3 题 6 分
%% typeId: 1
%% type: 单选题
%% chapter: 曲线运动
%% point: 平抛运动
%% method: 无
%% score: 6
%% degree: 550
%% duplicateId: 0
%% body:
如图，在摩托车越野赛途中的水平路段前方有一个坑，该坑沿摩托车前进方向的水平宽度为 $3h$，其左边缘 $a$ 点比右边缘 $b$ 点高 $0.5h$。若摩托车经过 $a$ 点时的动能为 $E_{1}$，它会落到坑内 $c$ 点。$c$ 与 $a$ 的水平距离和高度差均为 $h$；若经过 $a$ 点时的动能为 $E_{2}$，该摩托车恰能越过坑到达 $b$ 点。$\frac{E_{2}}{E_{1}}$ 等于\xzanswer{B}

\onepicture[w=8.00cm]{\includesvg{figs/03}}

\fourchoices[answer=B]
{20}
{18}
{9.0}
{3.0}
%% answer: B
%% memo:
\memoanswer{
由题意可知，当在 $a$ 点动能为 $E_{1}$ 时，有

$E_{1}=\frac{1}{2}mv_{1}^{2}$

根据平抛运动规律有

$h=\frac{1}{2}gt_{1}^{2}$，$h=v_{1}t_{1}$

当在 $a$ 点动能为 $E_{2}$ 时，有

$E_{2}=\frac{1}{2}mv_{2}^{2}$

根据平抛运动规律有

$\frac{1}{2}h=\frac{1}{2}gt_{2}^{2}$，$3h=v_{2}t_{2}$

联立以上各式可解得 $\frac{E_{2}}{E_{1}}=18$

故选 B。
}

\item
%% number: 4
%% paperName: 2020年全国理综Ⅱ第 4 题 6 分
%% typeId: 1
%% type: 单选题
%% chapter: 磁场
%% point: 带电粒子在磁场中的运动
%% method: 无
%% score: 6
%% degree: 650
%% duplicateId: 0
%% body:
CT 扫描是计算机 X 射线断层扫描技术的简称，CT 扫描机可用于对多种病情的探测。图~\ref{2020全国理综Ⅱ04a}是某种 CT 机主要部分的剖面图，其中 X 射线产生部分的示意图如图~\ref{2020全国理综Ⅱ04b}所示。图~\ref{2020全国理综Ⅱ04b}中 M、N 之间有一电子束的加速电场，虚线框内有匀强偏转磁场；经调节后电子束从静止开始沿带箭头的实线所示的方向前进，打到靶上，产生 X 射线（如图中带箭头的虚线所示）；将电子束打到靶上的点记为 P 点。则\xzanswer{D}

\twopicture[labelA=2020全国理综Ⅱ04a, labelB=2020全国理综Ⅱ04b, wA=7.00cm, wB=4.70cm]{figs/04a.png}{figs/04b.png}

\fourchoices[answer=D]
{M 处的电势高于 N 处的电势}
{增大 M、N 之间的加速电压可使 P 点左移}
{偏转磁场的方向垂直于纸面向外}
{增大偏转磁场磁感应强度的大小可使 P 点左移}
%% answer: D
%% memo:
\memoanswer{
A．由于电子带负电，要在 M、N 间加速则 M、N 间电场方向由 N 指向 M，根据沿着电场线方向电势逐渐降低可知 M 的电势低于 N 的电势，故 A 错误；

B．增大加速电压则根据 $eU=\frac{1}{2}mv^{2}$ 可知会增大到达偏转磁场的速度；又根据在偏转磁场中洛伦兹力提供向心力有 $evB=m\frac{v^{2}}{R}$，可得 $R=\frac{mv}{eB}$，可知会增大在偏转磁场中的偏转半径，由于磁场宽度相同，故根据几何关系可知会减小偏转的角度，故 P 点会右移，故 B 错误；

C．电子在偏转磁场中做圆周运动，向下偏转，根据左手定则可知磁场方向垂直于纸面向里，故 C 错误；

D．由 B 选项的分析可知，当其他条件不变时，增大偏转磁场磁感应强度会减小半径，从而增大偏转角度，使 P 点左移，故 D 正确。

故选 D。
}

\item
%% number: 5
%% paperName: 2020年全国理综Ⅱ第 5 题 6 分
%% typeId: 1
%% type: 单选题
%% chapter: 原子和原子核
%% point: 人工核反应
%% method: 无
%% score: 6
%% degree: 650
%% duplicateId: 0
%% body:
氘核 $\ce{^{2}_{1}H}$ 可通过一系列聚变反应释放能量，其总效果可用反应式

$\ce{6^{2}_{1}H -> 2^{4}_{2}He + 2^{1}_{1}H + 2^{1}_{0}n} + 43.15\UMeV$

表示。海水中富含氘，已知 1 $\Ukg$ 海水中含有的氘核约为 $1.0\times10^{22}$ 个，若全都发生聚变反应，其释放的能量与质量为 $M$ 的标准煤燃烧时释放的热量相等；已知 1 $\Ukg$ 标准煤燃烧释放的热量约为 $2.9\times10^{7}\UJ$，1 $\UMeV$ $=1.6\times10^{-13}\UJ$，则 $M$ 约为\xzanswer{C}

\fourchoices[answer=C]
{40 $\Ukg$}
{100 $\Ukg$}
{400 $\Ukg$}
{1 000 $\Ukg$}
%% answer: C
%% memo:
\memoanswer{
氘核 $\ce{^{2}_{1}H}$ 可通过一系列聚变反应释放能量，其总效果可用反应式

$\ce{6^{2}_{1}H -> 2^{4}_{2}He + 2^{1}_{1}H + 2^{1}_{0}n} + 43.15\UMeV$

表示，则平均每个氘核聚变释放的能量为

$\varepsilon=\frac{E}{6}=\frac{43.15}{6}\UMeV$

1 $\Ukg$ 海水中含有的氘核约为 $1.0\times10^{22}$ 个，可以放出的总能量为

$E_{0}=N\varepsilon$

由 $Q=mq$ 可得，要释放相同的热量，需要燃烧标准煤的质量

$m=\frac{Q}{q}=\frac{E_{0}}{q}\approx400\Ukg$

故选 C。
}

\item
%% number: 6
%% paperName: 2020年全国理综Ⅱ第 6 题 6 分
%% typeId: 2
%% type: 多选题
%% chapter: 电磁感应
%% point: 远距离输电
%% method: 无
%% score: 6
%% degree: 650
%% duplicateId: 0
%% body:
特高压输电可使输送中的电能损耗和电压损失大幅降低。我国已成功掌握并实际应用了特高压输电技术。假设从 A 处采用 550 $\UkV$ 的超高压向 B 处输电，输电线上损耗的电功率为 $\Delta P$，到达 B 处时电压下降了 $\Delta U$。在保持 A 处输送的电功率和输电线电阻都不变的条件下，改用 1 100 $\UkV$ 特高压输电，输电线上损耗的电功率变为 $\Delta P'$，到达 B 处时电压下降了 $\Delta U'$。不考虑其他因素的影响，则\xzanswer{AD}

\fourchoices[answer=AD]
{$\Delta P'=\frac{1}{4}\Delta P$}
{$\Delta P'=\frac{1}{2}\Delta P$}
{$\Delta U'=\frac{1}{4}\Delta U$}
{$\Delta U'=\frac{1}{2}\Delta U$}
%% answer: AD
%% memo:
\memoanswer{
输电线上损失的功率 $\Delta P=\left(\frac{P}{U}\right)^{2}\cdot r$，损失的电压 $\Delta U=\frac{P}{U}\cdot r$。当输送电压变为原来的 2 倍，损失的功率变为原来的 $\frac{1}{4}$，即 $\Delta P'=\frac{1}{4}\Delta P$；损失的电压变为原来的 $\frac{1}{2}$，即 $\Delta U'=\frac{1}{2}\Delta U$。故选 AD。
}

\item
%% number: 7
%% paperName: 2020年全国理综Ⅱ第 7 题 6 分
%% typeId: 1
%% type: 单选题
%% chapter: 电场
%% point: 电场线
%% method: 无
%% score: 6
%% degree: 550
%% duplicateId: 0
%% body:
如图，竖直面内一绝缘细圆环的上、下半圆分别均匀分布着等量异种电荷。$a$、$b$ 为圆环水平直径上的两个点，$c$、$d$ 为竖直直径上的两个点，它们与圆心的距离均相等。则\xzanswer{ABC}

\onepicture[w=4.30cm]{figs/07.png}

\fourchoices[answer=ABC]
{$a$、$b$ 两点的场强相等}
{$a$、$b$ 两点的电势相等}
{$c$、$d$ 两点的场强相等}
{$c$、$d$ 两点的电势相等}
%% answer: ABC
%% memo:
\memoanswer{
BD．如下图所示，为等量异种电荷周围空间的电场分布图。本题的带电圆环，可拆解成这样无数对等量异种电荷的电场，沿竖直直径平行放置。它们有共同的对称轴 PP$'$，PP$'$ 所在的水平面与每一条电场线都垂直，即为等势面，延伸到无限远处，电势为零。故在 PP$'$ 上的点电势为零，即 $\varphi_{a}=\varphi_{b}=0$；而从 M 点到 N 点，电势一直在降低，即 $\varphi_{c}>\varphi_{d}$，故 B 正确，D 错误。

\onepicture[w=3.30cm]{figs/07_an.png}

AC．上下两侧电场线分布对称，左右两侧电场线分布也对称，由电场的叠加原理可知 AC 正确。

故选 ABC。
}

\item
%% number: 8
%% paperName: 2020年全国理综Ⅱ第 8 题 6 分
%% typeId: 2
%% type: 多选题
%% chapter: 运动和力
%% point: 动量守恒定律
%% method: 无
%% score: 6
%% degree: 350
%% duplicateId: 0
%% body:
水平冰面上有一固定的竖直挡板，一滑冰运动员面对挡板静止在冰面上，他把一质量为 4.0 $\Ukg$ 的静止物块以大小为 5.0 $\Ums$ 的速度沿与挡板垂直的方向推向挡板，运动员获得退行速度；物块与挡板弹性碰撞，速度反向，追上运动员时，运动员又把物块推向挡板，使其再一次以大小为 5.0 $\Ums$ 的速度与挡板弹性碰撞。总共经过 8 次这样推物块后，运动员退行速度的大小大于 5.0 $\Ums$，反弹的物块不能再追上运动员。不计冰面的摩擦力，该运动员的质量可能为\xzanswer{BC}

\fourchoices[answer=BC]
{48 $\Ukg$}
{53 $\Ukg$}
{58 $\Ukg$}
{63 $\Ukg$}
%% answer: BC
%% memo:
\memoanswer{
设运动员和物块的质量分别为 $m$、$m_{0}$，规定运动员运动的方向为正方向，运动员开始时静止。第一次将物块推出后，运动员和物块的速度大小分别为 $v_{1}$、$v_{0}$，则根据动量守恒定律

$0=mv_{1}-m_{0}v_{0}$

解得 $v_{1}=\frac{m_{0}}{m}v_{0}$

物块与弹性挡板撞击后，运动方向与运动员同向，当运动员再次推出物块时

$mv_{1}+m_{0}v_{0}=mv_{2}-m_{0}v_{0}$

解得 $v_{2}=\frac{3m_{0}}{m}v_{0}$

第 3 次推出后

$mv_{2}+m_{0}v_{0}=mv_{3}-m_{0}v_{0}$

解得 $v_{3}=\frac{5m_{0}}{m}v_{0}$

依次类推，第 8 次推出后，运动员的速度 $v_{8}=\frac{15m_{0}}{m}v_{0}$

根据题意可知 $v_{8}=\frac{15m_{0}}{m}v_{0}>5\Ums$

解得 $m<60\Ukg$

第 7 次运动员的速度一定小于 $5\Ums$，则 $v_{7}=\frac{13m_{0}}{m}v_{0}<5\Ums$，解得 $m>52\Ukg$。

综上所述，运动员的质量满足 $52\Ukg<m<60\Ukg$，AD 错误，BC 正确。

故选 BC。
}

\item
%% number: 9
%% paperName: 2020年全国理综Ⅱ第 9 题 5 分
%% typeId: 4
%% type: 实验题
%% chapter: 运动和力
%% point: 连接体
%% method: 无
%% score: 5
%% degree: 550
%% duplicateId: 0
%% body:
一细绳跨过悬挂的定滑轮，两端分别系有小球 A 和 B，如图所示。一实验小组用此装置测量小球 B 运动的加速度。

令两小球静止，细绳拉紧，然后释放小球，测得小球 B 释放时的高度 $h_{0}=0.590\Um$，下降一段距离后的高度 $h=0.100\Um$；由 $h_{0}$ 下降至 $h$ 所用的时间 $T=0.730\Us$。由此求得小球 B 加速度的大小为 $a=$\tkanswer[1.2cm]{1.84} $\Umsq$（保留 3 位有效数字）。

从实验室提供的数据得知，小球 A、B 的质量分别为 100.0 $\Ug$ 和 150.0 $\Ug$，当地重力加速度大小为 $g=9.80\Umsq$。根据牛顿第二定律计算可得小球 B 加速度的大小为 $a'=$\tkanswer[1.2cm]{1.96} $\Umsq$（保留 3 位有效数字）。

可以看出，$a'$ 与 $a$ 有明显差异，除实验中的偶然误差外，写出一条可能产生这一结果的原因：\tkanswer[3cm]{滑轮的轴不光滑，绳和滑轮之间有摩擦（或滑轮有质量）}。

\onepicture[w=2.60cm]{\includesvg{figs/09}}
%% answer: 
%% memo:
\memoanswer{
（1）由题意可知小球下降过程中做匀加速直线运动，故根据运动学公式有

$h_{0}-h=\frac{1}{2}aT^{2}$

代入数据解得 $a=1.84\Umsq$；

（2）根据牛顿第二定律可知，对小球 A 有

$T-m_{A}g=m_{A}a'$

对小球 B 有

$m_{B}g-T=m_{B}a'$

带入已知数据解得 $a'=1.96\Umsq$；

（3）在实验中绳和滑轮之间有摩擦会造成实际计算值偏小。
}

\item
%% number: 10
%% paperName: 2020年全国理综Ⅱ第 10 题 10 分
%% typeId: 4
%% type: 实验题
%% chapter: 电路
%% point: 伏安法测电阻
%% method: 无
%% score: 10
%% degree: 450
%% duplicateId: 0
%% body:
某同学要研究一小灯泡 L（3.6 $\UV$，0.30 $\UA$）的伏安特性。所用器材有：电流表 A$_{1}$（量程 200 $\UmA$，内阻 $R_{g1}=10.0\UO$），电流表 A$_{2}$（量程 500 $\UmA$，内阻 $R_{g2}=1.0\UO$）、定值电阻 $R_{0}$（阻值 $R_{0}=10.0\UO$）、滑动变阻器 $R_{1}$（最大阻值 $10\UO$）、电源 $E$（电动势 4.5 $\UV$，内阻很小）、开关 S 和若干导线。该同学设计的电路如图~\ref{2020全国理综Ⅱ10a}所示。

\begin{enumerate}
\item
根据图~\ref{2020全国理综Ⅱ10a}，在图~\ref{2020全国理综Ⅱ10b}的实物图中画出连线\tkanswer[3cm]{}。

\twopicture[labelA=2020全国理综Ⅱ10a, labelB=2020全国理综Ⅱ10b, wA=4.80cm, wB=4.90cm]{figs/10a.png}{figs/10b.png}

\item
若 $I_{1}$、$I_{2}$ 分别为流过电流表 A$_{1}$ 和 A$_{2}$ 的电流，利用 $I_{1}$、$I_{2}$、$R_{g1}$ 和 $R_{0}$ 写出：小灯泡两端的电压 $U=$\tkanswer[3.5cm]{$I_{1}(R_{g1}+R_{0})$}，流过小灯泡的电流 $I=$\tkanswer[2.5cm]{$I_{2}-I_{1}$}。为保证小灯泡的安全，$I_{1}$ 不能超过\tkanswer[1.5cm]{180} $\UmA$。

\item
实验时，调节滑动变阻器，使开关闭合后两电流表的示数为零。逐次改变滑动变阻器滑片位置并读取相应的 $I_{1}$ 和 $I_{2}$。所得实验数据在下表中给出。

\begin{center}
\begin{tabular}{|c|cccccc|}
\hline
$I_{1}/\UmA$ & 32 & 55 & 85 & 125 & 144 & 173 \\
\hline
$I_{2}/\UmA$ & 171 & 229 & 299 & 379 & 424 & 470 \\
\hline
\end{tabular}
\end{center}

根据实验数据可算得，当 $I_{1}=173\UmA$ 时，灯丝电阻 $R=$\tkanswer[1.5cm]{11.6} $\UO$（保留 1 位小数）。

\item
如果用另一个电阻替代定值电阻 $R_{0}$，其他不变，为了能够测量完整的伏安特性曲线，所用电阻的阻值不能小于\tkanswer[1.5cm]{8.0} $\UO$（保留 1 位小数）。
\end{enumerate}
%% answer: 
\jdanswer{
\begin{enumerate}
\item
实物连线如答图所示。

\item
$I_{1}(R_{g1}+R_{0})$，$I_{2}-I_{1}$，180。

\item
11.6。

\item
8.0。
\end{enumerate}
}
%% memo:
\memoanswer{
（1）根据电路图连接实物图，注意各导线不要交叉，连接电表时注意电表的正、负接线柱，实物连线如答图所示。

\onepicture[w=5.40cm]{figs/10_an.png}

（2）由电路图可知，电流表 A$_{1}$ 与定值电阻 $R_{0}$ 串联充当电压表使用，所以小灯泡两端的电压 $U=I_{1}(R_{g1}+R_{0})$，流过小灯泡的电流 $I=I_{2}-I_{1}$。为保证小灯泡的安全，小灯泡两端的电压不能超过 3.6 $\UV$，故当小灯泡两端的电压为 3.6 $\UV$ 时，流过电流表 A$_{1}$ 的电流

$I=\frac{3.6\UV}{R_{g1}+R_{0}}=0.18\UA=180\UmA$。

（3）当 $I_{1}=173\UmA$ 时，$I_{2}=470\UmA$，小灯泡的电阻

$R=\frac{I_{1}(R_{g1}+R_{0})}{I_{2}-I_{1}}=\frac{173\times10^{-3}\times(10.0+10.0)}{(470-173)\times10^{-3}}\UO\approx11.6\UO$。

（4）电流表 A$_{1}$ 与定值电阻 $R_{0}$ 串联充当电压表使用，即流过电流表 A$_{1}$ 的电流为满偏电流 200 $\UmA$ 时，电流表 A$_{1}$ 与定值电阻两端的电压不能小于 3.6 $\UV$，则 $R_{0}'\geq\frac{3.6\UV}{200\times10^{-3}\UA}-10\UO=8.0\UO$。
}

\item
%% number: 11
%% paperName: 2020年全国理综Ⅱ第 11 题 12 分
%% typeId: 1
%% type: 单选题
%% chapter: 磁场
%% point: 带电粒子在磁场中的运动
%% method: 无
%% score: 12
%% degree: 450
%% duplicateId: 0
%% body:
如图，在 $0\leq x\leq h$，$-\infty<y<+\infty$ 区域中存在方向垂直于纸面的匀强磁场，磁感应强度 $B$ 的大小可调，方向不变。一质量为 $m$，电荷量为 $q$（$q>0$）的粒子以速度 $v_{0}$ 从磁场区域左侧沿 $x$ 轴进入磁场，不计重力。

\begin{enumerate}
\item
若粒子经磁场偏转后穿过 $y$ 轴正半轴离开磁场，分析说明磁场的方向，并求在这种情况下磁感应强度的最小值 $B_{m}$；
\item
如果磁感应强度大小为 $\frac{B_{m}}{2}$，粒子将通过虚线所示边界上的一点离开磁场。求粒子在该点的运动方向与 $x$ 轴正方向的夹角及该点到 $x$ 轴的距离。
\end{enumerate}

\onepicture[w=5.00cm]{figs/11.png}
%% answer: B
\jdanswer{
\begin{enumerate}
\item
磁场方向垂直于纸面向里；$B_{m}=\frac{mv_{0}}{qh}$。
\item
$\alpha=\frac{\pi}{6}$；$y=(2-\sqrt{3})h$。
\end{enumerate}
}
%% memo:
\memoanswer{
（1）由题意，粒子刚进入磁场时应受到方向向上的洛伦兹力，因此磁场方向垂直于纸面向里。设粒子在磁场中做圆周运动的半径为 $R$，根据洛伦兹力公式和圆周运动规律，有

$qv_{0}B=m\frac{v_{0}^{2}}{R}$\ccnum{①}

由此可得

$R=\frac{mv_{0}}{qB}$\ccnum{②}

粒子穿过 $y$ 轴正半轴离开磁场，其在磁场中做圆周运动的圆心在 $y$ 轴正半轴上，半径应满足 $R\leq h$\ccnum{③}。

由题意，当磁感应强度大小为 $B_{m}$ 时，粒子的运动半径最大，由此得

$B_{m}=\frac{mv_{0}}{qh}$\ccnum{④}

（2）若磁感应强度大小为 $\frac{B_{m}}{2}$，粒子做圆周运动的圆心仍在 $y$ 轴正半轴上，由\ccnum{②}\ccnum{④}式可得，此时圆弧半径为 $R'=2h$\ccnum{⑤}。

粒子会穿过图中 P 点离开磁场，运动轨迹如图所示。设粒子在 P 点的运动方向与 $x$ 轴正方向的夹角为 $\alpha$，由几何关系

$\sin\alpha=\frac{h}{2h}=\frac{1}{2}$\ccnum{⑥}

即 $\alpha=\frac{\pi}{6}$\ccnum{⑦}。

\onepicture[w=5.00cm]{figs/11_an.png}

由几何关系可得，P 点与 $x$ 轴的距离为

$y=2h(1-\cos\alpha)$\ccnum{⑧}

联立\ccnum{⑦}\ccnum{⑧}式得

$y=(2-\sqrt{3})h$\ccnum{⑨}。
}

\item
%% number: 12
%% paperName: 2020年全国理综Ⅱ第 12 题 20 分
%% typeId: 6
%% type: 计算题
%% chapter: 运动和力
%% point: 牛顿定律的应用
%% method: 无
%% score: 20
%% degree: 350
%% duplicateId: 0
%% body:
如图，一竖直圆管质量为 $M$，下端距水平地面的高度为 $H$，顶端塞有一质量为 $m$ 的小球。圆管由静止自由下落，与地面发生多次弹性碰撞，且每次碰撞时间均极短；在运动过程中，管始终保持竖直。已知 $M=4m$，球和管之间的滑动摩擦力大小为 $4mg$，$g$ 为重力加速度的大小，不计空气阻力。

\begin{enumerate}
\item
求管第一次与地面碰撞后的瞬间，管和球各自的加速度大小；
\item
管第一次落地弹起后，在上升过程中球没有从管中滑出，求管上升的最大高度；
\item
管第二次落地弹起的上升过程中，球仍没有从管中滑出，求圆管长度应满足的条件。
\end{enumerate}

\onepicture[w=3.60cm]{figs/12.png}
%% answer: 
\jdanswer{
\begin{enumerate}
\item
$a_{1}=2g$，$a_{2}=3g$。

\item
$H_{1}=\frac{13}{25}H$。

\item
$L\geq\frac{152}{125}H$。
\end{enumerate}
}
%% memo:
\memoanswer{
（1）管第一次落地弹起的瞬间，小球仍然向下运动。设此时管的加速度大小为 $a_{1}$，方向向下；球的加速度大小为 $a_{2}$，方向向上；球与管之间的摩擦力大小为 $f$，由牛顿第二定律有

$Ma_{1}=Mg+f$\ccnum{①}

$ma_{2}=f-mg$\ccnum{②}

联立\ccnum{①}\ccnum{②}式并代入题给数据，得

$a_{1}=2g$，$a_{2}=3g$\ccnum{③}。

（2）管第一次碰地前与球的速度大小相同。由运动学公式得，碰地前瞬间它们的速度大小均为

$v_{0}=\sqrt{2gH}$\ccnum{④}

方向均向下。管弹起的瞬间，管的速度反向，球的速度方向依然向下。

设自弹起时经过时间 $t_{1}$，管与小球的速度刚好相同。取向上为正方向，由运动学公式有

$v_{0}-a_{1}t_{1}=-v_{0}+a_{2}t_{1}$\ccnum{⑤}

联立\ccnum{③}\ccnum{④}\ccnum{⑤}式得

$t_{1}=\frac{2}{5}\sqrt{\frac{2H}{g}}$\ccnum{⑥}。

设此时管下端距地面的高度为 $h_{1}$，速度为 $v$。由运动学公式可得

$h_{1}=v_{0}t_{1}-\frac{1}{2}a_{1}t_{1}^{2}$\ccnum{⑦}

$v=v_{0}-a_{1}t_{1}$\ccnum{⑧}

可判断此时 $v>0$。此后，管与小球将以加速度 $g$ 减速上升 $h_{2}$，到达最高点。由运动学公式有

$h_{2}=\frac{v^{2}}{2g}$\ccnum{⑨}

设管第一次落地弹起后上升的最大高度为 $H_{1}$，则

$H_{1}=h_{1}+h_{2}$\ccnum{⑩}

联立\ccnum{③}\ccnum{④}\ccnum{⑥}\ccnum{⑦}\ccnum{⑧}\ccnum{⑨}\ccnum{⑩}式可得

$H_{1}=\frac{13}{25}H$\ccnum{⑪}。

（3）设第一次弹起过程中球相对管的位移为 $x_{1}$。在管开始下落到上升 $H_{1}$ 这一过程中，由动能定理有

$Mg(H-H_{1})+mg(H-H_{1}+x_{1})-4mgx_{1}=0$\ccnum{⑫}

联立\ccnum{⑪}\ccnum{⑫}式并代入题给数据得

$x_{1}=\frac{4}{5}H$\ccnum{⑬}

同理可推得，管与球从再次下落到第二次弹起至最高点的过程中，球与管的相对位移 $x_{2}$ 为 $x_{2}=\frac{4}{5}H_{1}$\ccnum{⑭}。

设圆管长度为 $L$。管第二次落地弹起后的上升过程中，球不会滑出管外的条件是 $x_{1}+x_{2}\leq L$\ccnum{⑮}，联立\ccnum{⑪}\ccnum{⑬}\ccnum{⑭}\ccnum{⑮}式，解得 $L$ 应满足的条件为

$L\geq\frac{152}{125}H$\ccnum{⑯}。
}

\item
%% number: 13
%% paperName: 2020年全国理综Ⅱ第 13 题 10 分
%% typeId: 6
%% type: 计算题
%% chapter: 气体性质
%% point: 玻意耳定律
%% method: 无
%% score: 10
%% degree: 550
%% duplicateId: 0
%% body:
【物理选修 3-3】

\begin{enumerate}
\item
（5 分）下列关于能量转换过程的叙述，违背热力学第一定律的有\tkanswer[1.2cm]{B}，不违背热力学第一定律、但违背热力学第二定律的有\tkanswer[1.2cm]{C}。（填正确答案标号）

（A）汽车通过燃烧汽油获得动力并向空气中散热

（B）冷水倒入保温杯后，冷水和杯子的温度都变得更低

（C）某新型热机工作时将从高温热源吸收的热量全部转化为功，而不产生其他影响

（D）冰箱的制冷机工作时从箱内低温环境中提取热量散发到温度较高的室内

\item
（10 分）潜水钟是一种水下救生设备，它是一个底部开口、上部封闭的容器，外形与钟相似。潜水钟在水下时其内部上方空间里存有空气，以满足潜水员水下避险的需要。为计算方便，将潜水钟简化为截面积为 $S$、高度为 $h$、开口向下的圆筒；工作母船将潜水钟由水面上方开口向下吊放至深度为 $H$ 的水下，如图所示。已知水的密度为 $\rho$，重力加速度大小为 $g$，大气压强为 $p_{0}$，$H\gg h$，忽略温度的变化和水密度随深度的变化。

\begin{enumerate}
\item
求进入圆筒内水的高度 $l$；
\item
保持 $H$ 不变，压入空气使筒内的水全部排出，求压入的空气在其压强为 $p_{0}$ 时的体积。
\end{enumerate}

\onepicture[w=4.20cm]{figs/13.png}
\end{enumerate}
%% answer: 
\jdanswer{
\begin{enumerate}
\item
B；C。

\item
$l=\frac{\rho gH}{p_{0}+\rho gH}h$，$V=\frac{\rho gSHh}{p_{0}}$。
\end{enumerate}
}
%% memo:
\memoanswer{
（1）A．燃烧汽油产生的内能一方面向机械能转化，同时热传递向空气转移，既不违背热力学第一定律，也不违背热力学第二定律；

B．冷水倒入保温杯后，没有对外做功，同时也没有热传递，内能不可能减少，故违背热力学第一定律；

C．某新型热机工作时将从高温热源吸收的热量全部转化为功，必然产生其他影响，故违背热力学第二定律；

D．制冷机消耗电能工作时从箱内低温环境中提取热量散发到温度较高的室内，发生了内能的转移，同时对外界产生了影响，既不违背热力学第一定律，也不违背热力学第二定律。

故违背热力学第一定律的有 B，不违背热力学第一定律、但违背热力学第二定律的有 C。

（2）设潜水钟在水面上方时和放入水下后筒内气体的体积分别为 $V_{0}$ 和 $V_{1}$，放入水下后筒内气体的压强为 $p_{1}$，由玻意耳定律和题给条件有

$p_{1}V_{1}=p_{0}V_{0}$\ccnum{①}

$V_{0}=hS$\ccnum{②}，$V_{1}=(h-l)S$\ccnum{③}

$p_{1}=p_{0}+\rho g(H-l)$\ccnum{④}

联立以上各式并考虑到 $H\gg h>l$，解得

$l=\frac{\rho gH}{p_{0}+\rho gH}h$\ccnum{⑤}

设水全部排出后筒内气体的压强为 $p_{2}$，此时筒内气体的体积为 $V_{0}$，这些气体在其压强为 $p_{0}$ 时的体积为 $V_{3}$，由玻意耳定律有

$p_{2}V_{0}=p_{0}V_{3}$\ccnum{⑥}

其中 $p_{2}=p_{0}+\rho gH$\ccnum{⑦}。

设需压入筒内的气体体积为 $V$，依题意

$V=V_{3}-V_{0}$\ccnum{⑧}

联立\ccnum{②}\ccnum{⑥}\ccnum{⑦}\ccnum{⑧}式得

$V=\frac{\rho gSHh}{p_{0}}$\ccnum{⑨}。
}

\item
%% number: 14
%% paperName: 2020年全国理综Ⅱ第 14 题 10 分
%% typeId: 6
%% type: 计算题
%% chapter: 光学
%% point: 全反射
%% method: 无
%% score: 10
%% degree: 550
%% duplicateId: 0
%% body:
【物理选修 3-4】

\begin{enumerate}
\item
（5 分）用一个摆长为 80.0 $\Ucm$ 的单摆做实验，要求摆动的最大角度小于 $5^\circ$，则开始时将摆球拉离平衡位置的距离应不超过\tkanswer[1.5cm]{6.9} $\Ucm$（保留 1 位小数）。（提示：单摆被拉开小角度的情况下，所求的距离约等于摆球沿圆弧移动的路程。）

某同学想设计一个新单摆，要求新单摆摆动 10 个周期的时间与原单摆摆动 11 个周期的时间相等。新单摆的摆长应该取为\tkanswer[1.5cm]{96.8} $\Ucm$。

\item
（10 分）直角棱镜的折射率 $n=1.5$，其横截面如图所示，图中 $\angle C=90^\circ$，$\angle A=30^\circ$。截面内一细束与 BC 边平行的光线，从棱镜 AB 边上的 D 点射入，经折射后射到 BC 边上。

\begin{enumerate}
\item
光线在 BC 边上是否会发生全反射？说明理由；
\item
不考虑多次反射，求从 AC 边射出的光线与最初的入射光线夹角的正弦值。
\end{enumerate}

\onepicture[w=4.60cm]{\includesvg{figs/14}}
\end{enumerate}
%% answer: 
\jdanswer{
\begin{enumerate}
\item
6.9，96.8。

\item
会发生全反射；$\sin r'=\frac{2\sqrt{2}-\sqrt{3}}{4}$。
\end{enumerate}
}
%% memo:
\memoanswer{
（1）开始时将摆球拉离平衡位置的距离最大，圆心角很小时，圆的弧长与弦长近似相等，则

$x=2\pi\times80\Ucm\times\frac{5^\circ}{360^\circ}\approx6.97\Ucm$

题中要求摆动的最大角度小于 $5^\circ$，且保留 1 位小数，所以拉离平衡位置的距离不超过 6.9 $\Ucm$；

根据单摆周期公式 $T=2\pi\sqrt{\frac{L}{g}}$ 结合题意可知 $10T'=11T$，代入数据为

$10\sqrt{L'}=11\sqrt{80\Ucm}$

解得新单摆的摆长为 $L'=96.8\Ucm$。

（2）如图，设光线在 D 点的入射角为 $i$，折射角为 $r$。折射光线射到 BC 边上的 E 点。设光线在 E 点的入射角为 $\theta$，由几何关系，有

$\theta=90^\circ-(30^\circ-r)>60^\circ$\ccnum{①}

根据题给数据得

$\sin\theta>\sin60^\circ>\frac{1}{n}$\ccnum{②}

即 $\theta$ 大于全反射临界角，因此光线在 E 点发生全反射。

\onepicture[w=8.00cm]{figs/14_an.png}

设光线在 AC 边上的 F 点射出棱镜，光线的入射角为 $i'$，折射角为 $r'$，由几何关系、反射定律及折射定律，有

$i=30^\circ$\ccnum{③}

$i'=90^\circ-\theta$\ccnum{④}

$\sin i=n\sin r$\ccnum{⑤}

$n\sin i'=\sin r'$\ccnum{⑥}

联立\ccnum{①}\ccnum{③}\ccnum{④}\ccnum{⑤}\ccnum{⑥}式并代入题给数据，得

$\sin r'=\frac{2\sqrt{2}-\sqrt{3}}{4}$\ccnum{⑦}

由几何关系知，$r'$ 即 AC 边射出的光线与最初的入射光线的夹角。
}

\end{enumerate}

\end{document}
